% ============================================================
% sec/06etica.tex
%
% OBLIGATORIO: riesgos identificados y forma prevista de
% mitigarlos. La detección corresponde al momento previo al
% entrenamiento.
% ============================================================
\section{Consideraciones Éticas Preliminares}
\label{sec:etica}

\subsection{Naturaleza de los datos}
\label{subsec:sensibilidad}

El conjunto de datos empleado procede de partidas entre agentes
automáticos, no entre personas. \pc{Confirmar que no contiene datos
personales identificables, y precisar la licencia bajo la cual fue
publicado y los usos que dicha licencia permite.}

\razon{La ausencia de datos personales debe declararse de forma explícita
en lugar de omitir el apartado. Una sección de ética que no se pronuncia
sobre la naturaleza de los datos deja abierta la pregunta.}

\subsection{Sesgo de la muestra}
\label{subsec:sesgo}

\pc{Criterio de selección aplicado por la fuente al publicar las
repeticiones.} En la medida en que dicho criterio privilegie a los agentes
mejor clasificados, la muestra no representa la población general de
participantes.

Saravanan y Guzdial orientan deliberadamente su agente a aproximar un nivel
de habilidad promedio, argumentando que el conjunto de estrategias
dominantes constituye una medida de popularidad determinada mayoritariamente
por jugadores de habilidad media \cite{saravanan_metadiscovery_2024}. Un
conjunto de datos restringido a la franja superior de clasificación observa,
en consecuencia, una población distinta de aquella que determina el
comportamiento agregado del juego.

Xenopoulos \textit{et al.} aportan evidencia sobre las consecuencias: la
importancia relativa de las características del estado difiere entre
poblaciones de distinto nivel de habilidad, y un modelo entrenado sobre una
de ellas no necesariamente se extiende a otra sin pérdida de desempeño
\cite{xenopoulos_winprob_2022}.

\subsection{Límites de generalización declarados}
\label{subsec:limitesgeneralizacion}

Los resultados de este trabajo no deben extrapolarse a:

\begin{itemize}
  \item Partidas entre personas. El conjunto procede de agentes
        automáticos, cuyo comportamiento no se asume equivalente al humano.
  \item Poblaciones de agentes de nivel de clasificación distinto del
        observado.
  \item Periodos posteriores a los cubiertos por el conjunto de datos, más
        allá de lo que el protocolo B permita evaluar.
  \item Versiones del juego o del simulador distintas de las presentes en
        los datos.
\end{itemize}

\razon{Este apartado existe para impedir que los resultados se empleen fuera
del ámbito en el cual fueron validados. Su omisión trasladaría al lector la
responsabilidad de inferir los límites.}

\subsection{Riesgos de calidad no aislables}
\label{subsec:calidadnoaislable}

\pc{Descripción de las anomalías detectadas en la auditoría que no admiten
aislamiento completo: episodios interrumpidos, comportamientos degenerados
de los agentes, o duraciones anómalas.}

Xenopoulos \textit{et al.} reconocen una limitación análoga al señalar que
no les resulta posible identificar a los jugadores que incurren en trampa,
más frecuentes en una de las poblaciones estudiadas
\cite{xenopoulos_winprob_2022}.

\razon{Declarar las anomalías que se sabe que existen pero no se pueden
aislar por completo es preferible a presentar el conjunto de datos como
libre de contaminación. El precedente citado establece que esta declaración
es práctica aceptada en el dominio.}

\subsection{Mitigación prevista antes del entrenamiento}
\label{subsec:mitigacionetica}

\begin{itemize}
  \item Declaración explícita de la población observada y de los límites de
        extrapolación, incorporada al artículo y a la documentación del
        repositorio.
  \item Evaluación bajo el protocolo B, que mide de forma directa la
        degradación del desempeño ante un cambio de la población observada.
  \item Reporte separado de la calidad de la calibración, dado que su
        deterioro bajo cambio de población puede ser considerablemente
        mayor que el de la capacidad de discriminación
        \cite{xenopoulos_winprob_2022}.
\end{itemize}

\razon{Las tres mitigaciones son verificables en el propio artículo. Una
mitigación cuya ejecución no puede comprobarse no constituye una
mitigación.}

\subsection{Aspectos diferidos}
\label{subsec:diferido}

La ética del despliegue, los escenarios de uso indebido y los riesgos
operativos asociados a la integración del modelo en un agente inteligente
se abordan en la etapa final del proyecto.